\chapter{Algèbre généalogique de chemins et représentation de Feynman}
\label{chap:integral-genealogica}

La formulation de Feynman attribue à chaque histoire possible une amplitude et somme
ces amplitudes pour obtenir le propagateur. En HMT, le mot \emph{histoire}
possède un contenu antérieur à l'amplitude : un chemin conserve position, feuillet,
orientation, phase, retenue, frontière, mémoire et condition de survie. La
somme de chemins ne s'ajoute pas de l'extérieur à ce système. Elle apparaît lorsqu'on représente
linéairement sa catégorie de chemins
\cite{Feynman1948,FeynmanHibbs1965}.

Le développement possède trois niveaux mathématiques liés. Le développement d'une
puissance matricielle comme somme finie de chemins est déjà démontré en mécanique
dimensionnelle. Le ledger \(27\times27\) fournit une réalisation unitaire
avec shift, mélangeur et phase, diagonalisable par Fourier. La famille de sommes
finies admet ensuite une limite cylindrique lorsque la mesure, l'action et le
raffinement satisfont les hypothèses du théorème de convergence. Ainsi se trouvent
réunies l'identité combinatoire, sa réalisation spectrale et son extension
fonctionnelle.

\section{La catégorie enrichie de chemins}

Soit \(X\) un ensemble fini d'états observables et soit
\(\mathsf P_{\TPK}\) la catégorie libre engendrée par les transitions
admissibles du TPK. Notons
\(\widetilde X=\operatorname{Ob}(\mathsf P_{\TPK})\). Ses objets sont des états
enrichis de frontière. Un morphisme
\(\gamma:\widetilde x\to\widetilde y\) est une séquence composable de
transitions enrichies.
La composition est la concaténation et l'identité est le chemin vide.

La projection visible
\[
 q_{\rm vis}:\mathsf P_{\TPK}\longrightarrow \mathsf P_X
\]
oublie une partie de la fibre. Deux morphismes différents peuvent avoir le même
mot observable. Cette non-fidélité est essentielle : une somme qui agrégerait prématurément
sur \(\mathsf P_X\) perdrait précisément l'information capable de
distinguer des futurs incompatibles.

\begin{definition}[Poids généalogique]
Soit \(\mathcal B\) une algèbre réelle normée, non nécessairement commutative, et
soit \(B(\mathcal B,\cdot,1)\) la catégorie à un objet dont le monoïde
d'endomorphismes est \((\mathcal B,\cdot,1)\). Un poids généalogique est un foncteur
\[
 W:\mathsf P_{\TPK}\longrightarrow B(\mathcal B,\cdot,1)
\]
qui satisfait
\[
 W(\id_x)=1_{\mathcal B},\qquad
 W(\gamma_2\circ\gamma_1)=W(\gamma_2)W(\gamma_1).
\]
L'ordre du produit suit l'ordre causal du chemin.
\end{definition}

Le poids peut contenir un transport de fibre, une pénalisation positive, une
action, un caractère du ledger ou plusieurs de ces facteurs. La définition ne
confond pas ces fonctions : chaque composante conserve son domaine et sa règle de
composition.

\section{La somme finie exacte}

Soit \(E_x\) la fibre linéaire sur \(x\in X\) et
\(\mathcal H_X=\bigoplus_{x\in X}E_x\). Soit
\(U=(U_{yx})_{x,y\in X}\) l'opérateur global par blocs, avec
\(U_{yx}:E_x\to E_y\). Nous interprétons \(U_{yx}\) comme le transport élémentaire
de \(x\) à \(y\). Pour un chemin visible
\(\gamma=(x_0,x_1,\ldots,x_R)\), nous définissons
\[
 W_U(\gamma)=U_{x_Rx_{R-1}}\cdots U_{x_1x_0}.
\]

\begin{theorem}[Développement fini de chemins]
\label{thm:suma-finita-caminos}
Pour tout \(R\geq 1\),
\begin{equation}
 (U^R)_{fi}
 =\sum_{\substack{\gamma:i\to f\\|\gamma|=R}}W_U(\gamma).
 \label{eq:suma-finita-caminos}
\end{equation}
Si l'opérateur global \(U:\mathcal H_X\to\mathcal H_X\) est orthogonal ---ou
unitaire après le choix d'une carte complexe---, l'évolution
\(\psi_R=U^R\psi_0\) conserve la norme.
\end{theorem}

\begin{proof}
L'entrée \((f,i)\) du produit \(U^R\) est
\[
 \sum_{x_1,\ldots,x_{R-1}}
 U_{fx_{R-1}}U_{x_{R-1}x_{R-2}}\cdots U_{x_1i}.
\]
Chaque choix d'indices intermédiaires détermine un unique chemin de longueur
\(R\), et chaque chemin détermine un unique monôme ordonné. La conservation de la
norme est une conséquence immédiate de \((U^R)^{\mathsf T}U^R=I\), ou de sa version
unitaire.
\end{proof}

La formule \eqref{eq:suma-finita-caminos} est l'identité combinatoire qui
sous-tend la somme d'histoires. Son contenu HMT apparaît lorsqu'on raffine chaque
chemin visible par les fibres de \(q_{\rm vis}\) :
\begin{equation}
 \mathcal K_R(\widetilde f,\widetilde i)
 =\sum_{\substack{\gamma\in
        \mathsf P_{\TPK}(\widetilde i,\widetilde f)\\|\gamma|=R}}
 W(\gamma).
 \label{eq:kernel-genealogico-finito}
\end{equation}
Dans cette somme, deux chemins publiés avec les mêmes extrémités restent des termes
distincts tant qu'ils diffèrent par leur généalogie.

\section{Réalisation unitaire sur le registre généalogique \(27\times27\)}

La source historique développe une marche quantique sur le tore
\[
 \Lambda_{27}=(\ZZ/27\ZZ)^2
\]
avec quatre orientations cardinales
\(\mathcal D=\{N,S,E,W\}\). Son espace réel de phase est
\[
 \mathcal H_{27}^{\RR}
 =\ell^2(\Lambda_{27})\otimes\RR^4\otimes E,
\]
où \((E,J_E)\) réalise intérieurement la structure complexe. Le mélangeur
de directions est la matrice de Walsh--Hadamard
\[
 C_4=\frac12
 \begin{pmatrix}
 1&1&1&1\\
 1&1&-1&-1\\
 1&-1&1&-1\\
 1&-1&-1&1
 \end{pmatrix},
 \qquad C_4^{\mathsf T}C_4=I_4.
\]
Une action élémentaire orientée \(s_d\in\mathcal L_S\) définit la porte
réelle de phase
\[
 D_J=\operatorname{diag}
 \left(
  e^{-J_Es_N/\hbar_{\rm int}},
  e^{-J_Es_S/\hbar_{\rm int}},
  e^{-J_Es_E/\hbar_{\rm int}},
  e^{-J_Es_W/\hbar_{\rm int}}
 \right).
\]
Le shift conditionné transporte chaque composante dans sa direction :
\begin{align*}
 (S\psi)(x,y,N)&=\psi(x,y-1,N),&
 (S\psi)(x,y,S)&=\psi(x,y+1,S),\\
 (S\psi)(x,y,E)&=\psi(x-1,y,E),&
 (S\psi)(x,y,W)&=\psi(x+1,y,W),
\end{align*}
avec des coordonnées modulo \(27\). Le pas élémentaire et la composition
dodécaphasique sont
\begin{equation}
 U_{\rm micro}=S D_J C_4,
 \qquad
 U_{12}=U_{\rm micro}^{12}.
 \label{eq:unitario-ledger-27}
\end{equation}

\begin{proposition}[Unitarité et somme exacte dans le ledger]
L'opérateur \(U_{\rm micro}\) est orthogonal sur la réalification et unitaire
dans la carte complexe. Pour tout \(n\), chaque entrée de
\(U_{\rm micro}^n\) est la somme ordonnée des poids de tous les chemins de
\(n\) pas qui relient ses états extrêmes.
\end{proposition}

\begin{proof}
Le shift est une permutation, \(C_4\) est orthogonal et chaque bloc
\(e^{-J_Es_d/\hbar_{\rm int}}\) est une rotation. Leur produit est orthogonal.
La deuxième affirmation est le
\cref{thm:suma-finita-caminos} appliqué à la matrice par blocs
\(U_{\rm micro}\).
\end{proof}

La symétrie translationnelle réduit le problème spectral à des blocs de taille
quatre. Pour
\[
 k_x=\frac{2\pi m}{27},\qquad
 k_y=\frac{2\pi n}{27},
 \qquad 0\le m,n<27,
\]
le shift se diagonalise comme
\[
 S(k_x,k_y)=\operatorname{diag}
 \bigl(e^{J_Ek_y},e^{-J_Ek_y},e^{J_Ek_x},e^{-J_Ek_x}\bigr),
\]
et
\begin{equation}
 U_{\rm micro}(k_x,k_y)=S(k_x,k_y)D_JC_4,
 \qquad
 U_{12}(k_x,k_y)=U_{\rm micro}(k_x,k_y)^{12}.
 \label{eq:bloque-fourier-ledger-27}
\end{equation}
Le spectre global est l'union des spectres des \(27^2\) blocs.
L'expression combine deux représentations de la même dynamique : la puissance
matricielle énumère les chemins et la transformée de Fourier publie les
quasi-énergies.

Dans chaque bloc unitaire, on peut choisir un logarithme spectral et définir
\[
 H_{\rm eff}(k_x,k_y)
 =\frac{\hbar_{\rm int}}{J_E\,\Delta t}
   \log U_{12}(k_x,k_y),
\]
compris au moyen du calcul fonctionnel dans la carte complexe ou sa réalification.
Le choix de branche fixe la zone de quasi-énergie. Le Hamiltonien effectif
émerge ainsi du propagateur déjà construit ; son exponentielle restitue
\(U_{12}\) dans chaque mode.

\section{Caractère réel d'action}

Soit \(\mathcal A_*:\mathsf P_{\TPK}\to\RR_{\geq0}\) le cocycle additif
d'action de la source HMT--MD. La jauge est
\[
 \lambda_*
 =1-\frac5{468}-\frac3{11}\Delta_4
 =0.9892859130191415\ldots\in(0,1).
\]
La clause locale \(\mathsf H_{\rm loc}\) attribue un coût unitaire à une rotation,
un coût \(\lambda_*\) à un changement de feuillet et additionne les deux lorsqu'ils se produisent sur la
même arête. Sous cette clause,
\[
 \mathcal A_*(\gamma)
 =\#\operatorname{turns}(\gamma)
  +\lambda_*\#\operatorname{sheet\mbox{-}switches}(\gamma),
\qquad
 \mathcal A_*(\gamma_2\circ\gamma_1)
 =\mathcal A_*(\gamma_2)+\mathcal A_*(\gamma_1).
\]
L'action dimensionnelle associée est
\[
 S_{\rm int}(\gamma)=\hbar_{\rm int}\mathcal A_*(\gamma).
\]

Dans le module de phase réelle \((E,J_E)\) du chapitre précédent, avec \(E\) un
espace euclidien réel de dimension finie, nous définissons
\begin{equation}
 \Phi_J(\gamma)
 :=\exp\!\left(-J_E\frac{S_{\rm int}(\gamma)}{\hbar_{\rm int}}\right)
 =\cos\mathcal A_*(\gamma)\,I
  -\sin\mathcal A_*(\gamma)\,J_E.
 \label{eq:caracter-real-accion}
\end{equation}

\begin{proposition}[Multiplicativité de la phase interne]
\label{prop:multiplicatividad-fase-interna}
L'application \(\Phi_J\) est un foncteur de \(\mathsf P_{\TPK}\) vers le groupe
orthogonal engendré par \(J_E\) :
\[
 \Phi_J(\id_x)=I,
 \qquad
 \Phi_J(\gamma_2\circ\gamma_1)
 =\Phi_J(\gamma_2)\Phi_J(\gamma_1).
\]
\end{proposition}

\begin{proof}
La première égalité utilise \(\mathcal A_*(\id_x)=0\). Pour la deuxième,
l'additivité de \(\mathcal A_*\) et la commutation des multiples du même
endomorphisme \(J_E\) donnent
\[
 e^{-J_E(\mathcal A_*(\gamma_2)+\mathcal A_*(\gamma_1))}
 =e^{-J_E\mathcal A_*(\gamma_2)}
  e^{-J_E\mathcal A_*(\gamma_1)}.
\]
\end{proof}

Soit maintenant \(T(e)\in\End(E)\) le transport borné de fibre d'une arête et
supposons
qu'il préserve \(J_E\). Le poids
\[
 W_J(\gamma)
 =\prod_{e\in\gamma}^{\longleftarrow}
   T(e)\exp\!\left(-J_E\frac{s(e)}{\hbar_{\rm int}}\right)
\]
réalise entièrement sur \(E\) la convention hamiltonienne
\(\exp(-iS_H[\gamma]/\hbar)\). La convention lagrangienne habituelle de Feynman
\(\exp(+iS_L[\gamma]/\hbar)\) s'obtient en prenant \(S_L=-S_H\), ou,
de manière équivalente, en introduisant une orientation
\(\varepsilon\in\{-1,+1\}\) dans
\(\exp(\varepsilon J_ES/\hbar_{\rm int})\). Le symbole scalaire \(i\) est la
carte complexe de l'opération algébrique réalisée par \(J_E\). Le coût
positif \(\mathcal A_*\) devient une action physique lorsqu'on applique l'application de
réalisation correspondante.

\section{Composition de propagateurs sur des chemins concaténés et loi des amplitudes}

La généalogie conserve la loi de composition qui rend le noyau dynamique.

\begin{proposition}[Convolution et coupure temporelle]
\label{prop:chapman-kolmogorov-genealogico}
Supposons que \(\widetilde X\) est fini ; plus généralement, que la famille de
produits suivante est absolument sommable. Si tout chemin de longueur
\(R+S\) possède une unique coupure après \(R\) pas, alors
\begin{equation}
 \mathcal K_{R+S}(\widetilde f,\widetilde i)
 =\sum_{\widetilde x\in\widetilde X}
   \mathcal K_S(\widetilde f,\widetilde x)
   \mathcal K_R(\widetilde x,\widetilde i).
 \label{eq:convolucion-kernel}
\end{equation}
L'identité reste valable avec des poids non commutatifs si l'on conserve
l'ordre causal des facteurs.
\end{proposition}

\begin{proof}
Nous partitionnons l'ensemble des chemins selon l'état enrichi
\(\widetilde x\) atteint à la coupure.
La factorisation fonctorielle du poids transforme la somme sur les chemins complets
en produit des deux sommes partielles, avec le segment postérieur à
gauche.
\end{proof}

La coupure sur l'état visible ne serait valable que si le poids et la composition
descendaient fidèlement par \(q_{\rm vis}\). En présence de mémoire, deux
tronçons ayant le même visible peuvent appartenir à des fibres incompatibles ; c'est pourquoi
la convolution s'effectue sur \(\widetilde X\).

Cette équation distingue la somme de chemins d'une collection de coïncidences
numériques. Le propagateur possède un semi-groupe, des extrémités, une composition et une loi de
raffinement.

\section{Limite cylindrique du système fini et construction de l'intégrale correspondante}

Une intégrale fonctionnelle exige de contrôler l'espace des chemins et son
raffinement. Pour des extrémités enrichies \(\widetilde i,\widetilde f\), soient
\(\mathcal P_n(\widetilde i,\widetilde f)\) les ensembles finis de chemins
admissibles à la résolution \(n\), avec des restrictions surjectives, et définissons
\[
 \Omega_{fi}^{\rm path}
 :=\varprojlim_n\mathcal P_n(\widetilde i,\widetilde f).
\]
C'est l'espace compact des généalogies de chemins ayant ces données de frontière.
Soit \(\{\Pi_n\}_{n\geq0}\) une suite de partitions cylindriques de
\(\Omega_{fi}^{\rm path}\), où \(\Pi_{n+1}\) raffine \(\Pi_n\). Pour chaque
cylindre \(C\in\Pi_n\), choisissons un représentant
\(\gamma_C\), un poids projectif \(\mu_n(C)\) et une action approchée
\(S_n(C)\).

Nous définissons l'opérateur simple
\begin{equation}
 \mathcal I_n(F)
 :=\sum_{C\in\Pi_n}
   \mu_n(C)
   F(\gamma_C)
   \exp\!\left(-J_E\frac{S_n(C)}{\hbar_{\rm int}}\right).
 \label{eq:integral-cilindrica-n}
\end{equation}

\begin{theorem}[Limite généalogique d'intégrales cylindriques]
\label{thm:integral-cilindrica-limite}
Supposons que :
\begin{enumerate}
 \item les masses cylindriques sont projectivement compatibles ;
 \item \(F:\Omega_{fi}^{\rm path}\to\End(E)\) est continue et bornée ;
 \item il existe une fonction continue
       \(S:\Omega_{fi}^{\rm path}\to\mathcal L_S\) telle que
       \[
        \max_{C\in\Pi_n}\sup_{\gamma\in C}
        \left\|
          \frac{S_n(C)-S(\gamma)}{\hbar_{\rm int}}
        \right\|\longrightarrow0;
       \]
 \item le pas des partitions,
       \(\operatorname{mesh}(\Pi_n)
        =\max_{C\in\Pi_n}\diam_{d_\vartheta}C\),
       tend vers zéro.
\end{enumerate}
Comme \(E\) est de dimension finie, \(\End(E)\), muni de la norme
d'opérateur, est un espace de Banach ; l'intégrale qui suit est donc une
intégrale de Bochner bien définie \cite{Bochner1955}.
Alors \(\mathcal I_n(F)\) converge en norme d'opérateur et sa limite ne
dépend pas des représentants. Nous écrivons
\begin{equation}
 \int_{\Omega_{fi}^{\rm path}}
 F(\gamma)
 \exp\!\left(-J_E\frac{S(\gamma)}{\hbar_{\rm int}}\right)
 \dd\mu(\gamma)
 :=\lim_{n\to\infty}\mathcal I_n(F).
 \label{eq:integral-genealogica}
\end{equation}
\end{theorem}

\begin{proof}
La compatibilité projective construit la mesure de cylindres \(\mu\). La
fonction
\[
 G(\gamma)=F(\gamma)
 \exp\!\left(-J_E S(\gamma)/\hbar_{\rm int}\right)
\]
est continue sur le compact \(\Omega_{fi}^{\rm path}\), donc uniformément
continue et bornée. Soit
\[
 G_n(C)=F(\gamma_C)
 \exp\!\left(-J_E S_n(C)/\hbar_{\rm int}\right).
\]
Comme \(J_E\) est une structure orthogonale,
\[
 \|e^{-J_Ea}-e^{-J_Eb}\|\leq |a-b|.
\]
La troisième hypothèse implique que
\(G_n(C)-G(\gamma_C)\) tend uniformément vers zéro. La quatrième hypothèse et la
continuité uniforme de \(G\) font des sommes avec \(G(\gamma_C)\) des
sommes de Riemann convergentes et indépendantes des représentants. La
limite coïncide avec l'intégrale de Bochner de \(G\) par rapport à \(\mu\).
Plus précisément, si
\[
 \epsilon_n
 =\max_{C\in\Pi_n}\sup_{\gamma\in C}
   \left\|
    \frac{S_n(C)-S(\gamma)}{\hbar_{\rm int}}
   \right\|,
\]
alors
\[
 \left\|\mathcal I_n(F)-\int G\,\dd\mu\right\|
 \leq
 \omega_F\!\bigl(\operatorname{mesh}(\Pi_n)\bigr)
 +\|F\|_\infty\epsilon_n,
\]
car l'intégrale se décompose sur les cylindres et la masse totale vaut un.
Les deux termes tendent vers zéro par hypothèse.
\end{proof}

\begin{corollary}[Convergence d'une famille de noyaux discrets]
\label{cor:convergencia-kernels-discretos}
Outre les hypothèses du
\cref{thm:integral-cilindrica-limite}, supposons qu'il existe une suite
cofinale de profondeurs \(R_n\) et des fonctions simples \(F_{fi,n}\) telles que
\[
 \mathcal K_{R_n}(\widetilde f,\widetilde i)
 =\mathcal I_n(F_{fi,n}),
 \qquad
 F_{fi,n}\longrightarrow F_{fi}
\]
uniformément. Alors les noyaux
\(\mathcal K_{R_n}(\widetilde f,\widetilde i)\) convergent vers l'intégrale
généalogique de \(F_{fi}\).
\end{corollary}

\begin{proof}
La différence entre \(\mathcal I_n(F_{fi,n})\) et
\(\mathcal I_n(F_{fi})\) est majorée par
\(\|F_{fi,n}-F_{fi}\|_\infty\), car \(\mu_n\) a une masse totale égale à un et la
phase est orthogonale. Le deuxième terme converge par le
\cref{thm:integral-cilindrica-limite}.
\end{proof}

Le théorème produit une intégrale rigoureuse sur l'espace généalogique
compact. Le corollaire explicite le pont qui identifie chaque noyau
matriciel à une somme cylindrique normalisée. L'intégrale oscillatoire de
Feynman en temps réel apparaît lorsqu'on choisit
une représentation physique et un raffinement dont l'action discrète approche
l'action du système. Cette identification exige de démontrer la convergence du
propagateur correspondant ; elle n'est pas accordée par la simple existence de
\(\mu\).

\section{Intégrale visible et intégrale conditionnée par la mémoire}

Un lecteur réel de chemins
\(\Val_{\rm path}:\Omega_{fi}^{\rm path}\to K\subset\RR\) induit la mesure
visible \(\nu=(\Val_{\rm path})_*\mu\). Pour une fonction qui se factorise par la
valeur, \(F=\widehat F\circ\Val_{\rm path}\),
\[
 \int_{\Omega_{fi}^{\rm path}}F\,\dd\mu
 =\int_K\widehat F(x)\,\dd\nu(x).
\]
Cependant, la phase ou le transport peuvent ne pas se factoriser par
\(\Val_{\rm path}\). La
désintégration de \(\mu\) sur les fibres de la valeur exprime la différence :
\begin{equation}
 \int_{\Omega_{fi}^{\rm path}}G(\gamma)\,\dd\mu(\gamma)
 =\int_K
   \left[
    \int_{\Val_{\rm path}^{-1}(x)}G(\gamma)\,\dd\mu_x(\gamma)
   \right]\dd\nu(x).
 \label{eq:desintegracion-caminos}
\end{equation}
L'intégrale intérieure conserve les histoires qui publient la même coordonnée.
C'est la forme mathématique de l'intégration par rapport aux valeurs sans réduire le
nombre à son chiffre.

\section{Action, interférence et persistance}

La somme généalogique conserve simultanément trois opérations :
\[
 \text{admissibilité de route}
 \quad\longrightarrow\quad
 \text{poids/action}
 \quad\longrightarrow\quad
 \text{interférence dans la représentation}.
\]
La première appartient au TPK et à sa loi de survie. La deuxième appartient
au cocycle et au transport dimensionnel. La troisième appartient au module de
phase et au lecteur physique. Aucune grandeur observée ne sélectionne
rétrospectivement le chemin qui la produit.

\begin{resultbox}
La somme finie de Feynman s'obtient comme représentation linéaire exacte de la
catégorie enrichie de chemins. La structure complexe peut être réalisée par un
endomorphisme réel \(J_E^2=-I\). Sous compatibilité des mesures, contraction des
cylindres et convergence uniforme des actions, les sommes cylindriques définissent
une intégrale généalogique. Les noyaux finis y convergent lorsqu'ils
satisfont, en outre, le pont du
\cref{cor:convergencia-kernels-discretos}.
\end{resultbox}

\begin{statusbox}
Le développement matriciel fini, le cocycle d'action et les réalisations réelles
de la structure complexe sont des résultats récupérés. Leur assemblage dans
\eqref{eq:caracter-real-accion} et \eqref{eq:kernel-genealogico-finito} est une
nouvelle formalisation exacte. Le théorème \ref{thm:integral-cilindrica-limite} est
exact sous ses prémisses. L'équivalence avec un modèle quantique continu
concret exige de vérifier ces prémisses pour le Hamiltonien et l'action de
ce modèle.
\end{statusbox}

Les prémisses du théorème identifient exactement la structure qui passe du
ledger fini à l'espace fonctionnel : compatibilité projective des poids,
contraction du pas, convergence uniforme de l'action et correspondance
entre noyaux et sommes cylindriques. Dans la réalisation
\eqref{eq:unitario-ledger-27}, la composition et l'unitarité sont des identités
algébriques ; en raffinant le ledger, les autres prémisses déterminent la
convergence vers le propagateur continu choisi.
